% Adaptador único del capítulo 38 del sucesor 102.
% Sustituye exclusivamente la jerarquía de encabezados. Los cuerpos de los
% dos propietarios científicos y del delta Ley 9 se conservan línea a línea,
% reordenados conforme al DAG causal aprobado. Véase MAPA_SEGMENTOS_CAPITULO_38.tsv.

% HMT_SOURCE_SEGMENT_BEGIN id=B00 source=colaboracion/parte_iii/source/public_final/base_83/c30_body.tex body_lines=1-2
\label{chap:p3-83-30-angulo_contraangulo}

% HMT_SOURCE_SEGMENT_END id=B00
% HMT_SOURCE_SEGMENT_BEGIN id=A00 source=colaboracion/parte_iii/source/public_final/ascensos_20260826/16f_realizacion_analitica_contraangulo.tex body_lines=1-4
\label{p3-20260826:chap:realizacion-analitica-contraangulo}
\index{contre-angle}
\index{microfibre de pi@microfibre de \(\pi\)!caractère analytique}
\index{caractère déterminantal}
% HMT_SOURCE_SEGMENT_END id=A00
% HMT_SOURCE_SEGMENT_BEGIN id=D00 source=manuscrito/sucesor_102/deltas_ley9/c38_par_angular_cartas_y_canales.tex body_lines=1
% Distribución causal exacta del delta Ley 9; contenido conservado sin síntesis.
% HMT_SOURCE_SEGMENT_END id=D00

\section{Phase angulaire intrinsèque \texorpdfstring{\(\alpha_{\mathrm{HMT}}\)}{alpha HMT}, relèvement nonadique \texorpdfstring{\(10^3\)}{10^3} et précurseur \texorpdfstring{\(x_A\)}{x_A}}
\label{sec:s102:c38:1}

\subsection{Construction du précurseur analytique \texorpdfstring{\(x_A\)}{x_A}, de la contraction \texorpdfstring{\(D_A\)}{D_A} et de la correction régionale \texorpdfstring{\(\mathscr C_\pi(\alpha)\)}{C_pi(alpha)}}
\label{sec:s102:c38:1:1}

% HMT_SOURCE_SEGMENT_BEGIN id=B01 source=colaboracion/parte_iii/source/public_final/base_83/c30_body.tex body_lines=4-32
\label{p3-83-c30-s1:chap:realizacion-analitica-contraangulo}
\index{coordonnée angulaire conjuguée \(C^*\)}
\index{microfibre de pi@microfibre de \(\pi\)!caractère analytique}
\index{caractère déterminantal}

Le catalogue fini et une évaluation analytique répondent à des questions de types
distincts. Le catalogue détermine quelles régions existent, lesquelles partagent un
retour et quelle information chacune conserve. Une évaluation analytique doit
en outre décider comment la longueur d'un mot se transforme en degré, comment
une branche se prolonge après son retour et avec quel caractère scalaire cette
prolongation se lit. Ce chapitre construit cette seconde opération de manière
explicite.

La construction part de quatre données déjà obtenues dans les chapitres
précédents : la microfibre à cinq régions de \(\pi\), la longueur six de
chaque mot régional, la clôture dodécaphasique de \(\alpha\) et le transport
bicouche associé à la coordonnée commune
\(A=10^3\alpha\). Aucune valeur de la coordonnée angulaire conjuguée \(C_*^{\rm HMT}\), de la composante réduite
d'action, de la fonctionnelle de Barbero--Immirzi ou d'une constante métrologique n'est
employée dans l'évaluation. Le résultat est un caractère régional convergent,
une récurrence exacte et une phase orientée dont la carte sexagésimale reproduit
la coordonnée angulaire conjuguée publiée.

La séparation entre la partie finie et la partie analytique est essentielle. L'entier
\(169\) sera un théorème du catalogue. La série qui le prolonge sera le
théorème d'un lecteur analytique dont les règles sont déclarées avant de l'évaluer.
Ainsi, les prémisses supplémentaires ne restent pas dissimulées dans une coïncidence
décimale.

% HMT_SOURCE_SEGMENT_END id=B01
\section{Construction préangulaire de \texorpdfstring{\(\eta_{\mathrm{ret}}\)}{eta_ret} au moyen de \texorpdfstring{\(H_5\)}{H5}, \texorpdfstring{\(D_A\)}{D_A} et \texorpdfstring{\(\mathscr C_\pi(\alpha)\)}{C_pi(alpha)}}
\label{sec:s102:c38:2}

\subsection{Dérivation indépendante de \texorpdfstring{\(\eta_{\mathrm{ret}}=H_5-(90/\pi)\mathscr C_\pi(\alpha)\)}{eta_ret=H5-(90/pi)C_pi(alpha)} au moyen du lecteur analytique régional}
\label{sec:s102:c38:2:1}

% HMT_SOURCE_SEGMENT_BEGIN id=A01 source=colaboracion/parte_iii/source/public_final/ascensos_20260826/16f_realizacion_analitica_contraangulo.tex body_lines=5-38,40-47

Le catalogue fini et une évaluation analytique répondent à des questions de types
distincts. Le catalogue détermine quelles régions existent, lesquelles partagent un
retour et quelle information chacune conserve. Une évaluation analytique doit
décider, en outre, comment la longueur d'un mot se transforme en degré, comment
une branche se prolonge après son retour et avec quel caractère scalaire on lit
cette prolongation. Cette section construit cette seconde opération de manière
explicite.

La construction part de la microfibre de cinq régions de \(\pi\), de la
longueur six de chaque mot régional, de la clôture dodécaphasique de \(\alpha\)
et de l'insertion angulaire déclarée
\[
 \iota_\angle(x):=10^3x\,u_\circ,\qquad
 \Adeg:=10^3\alpha .
\]
Ici, \(u_\circ\) est la base sexagésimale dont le tour complet a pour
coordonnée \(360\).
Le facteur \(10^3\) appartient à cette normalisation de carte et ne se déduit pas de la
clôture arithmétique de \(\alpha\). Le transport bicouche est construit ensuite
à partir de \(\Adeg\). La fonction qui évalue le lecteur ne reçoit comme
argument ni le contre-angle, ni la coordonnée postérieure au retour, ni la fonctionnelle
hexade--octade, ni une constante métrologique. Le résultat est un caractère
régional convergent, une récurrence exacte et une phase orientée dans la carte
sexagésimale déclarée.

La séparation entre la partie finie et la partie analytique est essentielle. L'entier
\(169\) est un théorème du catalogue. La série qui le prolonge appartient
à un lecteur analytique dont les règles sont déclarées avant son évaluation.

La phase est notée \(y_*\) et, lorsque l'on souligne son origine régionale,
\(y_{\rm reg}\) ; les deux symboles désignent le même lecteur construit
ci-après.


L'évaluation reçoit le catalogue de \(468\) émissions, la valeur exacte de
\(\alpha_{\rm HMT}\), \(\varphi=(1+\sqrt5)/2\), la formule fixée de \(H_5\),
la carte angulaire et les cinq règles du lecteur régional. Avec ces données,
elle calcule \(169\), \(D_A\), la récurrence et \(C^*\). Le théorème est exact dans
cette classe déclarée de lecteurs ; il n'utilise pas \(C^*\) comme entrée et n'affirme pas que
les règles de \(H_5\) constituent l'unique classe naturelle possible.

% HMT_SOURCE_SEGMENT_END id=A01
\subsection{Persistance du retour dans les \texorpdfstring{\(5\)}{5} régions de la microfibre}
\label{sec:s102:c38:2:2}

% HMT_SOURCE_SEGMENT_BEGIN id=B02 source=colaboracion/parte_iii/source/public_final/base_83/c30_body.tex body_lines=34-155

Soit
\[
\mathcal F_\pi=\Psi^{-1}(010211)\subset\mathcal U_6
\]
la microfibre relative au catalogue. Rappelons ses cinq éléments, maintenant séparés en bloc initial et bloc terminal :
\[
\begin{aligned}
 &(501,614,498\mid169,272,272),\\
 &(810,923,870\mid169,272,272),\\
 &(810,983,810\mid169,272,272),\\
 &(870,923,810\mid169,272,272),\\
 &(870,923,810\mid418,674,521).
\end{aligned}
\]
Nous définissons la projection de retour
\[
p_{\rm ret}:\mathbb Z^6\longrightarrow\mathbb Z^3,
\qquad
p_{\rm ret}(u_1,\ldots,u_6)=(u_4,u_5,u_6).
\]

\begin{proposition}[Bloc terminal persistant]
\label{p3-83-c30-s1:prop:bloque-terminal-persistente}
Le multiensemble \(p_{\rm ret}(\mathcal F_\pi)\) possède un unique élément de
multiplicité maximale :
\[
b_\pi=(169,272,272),
\qquad
\operatorname{mult}_{\mathcal F_\pi}(b_\pi)=4.
\]
L'unique bloc terminal restant est
\[
b_\pi^-=(418,674,521)
\]
et a une multiplicité de un.
\end{proposition}

\begin{proof}
L'affirmation s'obtient par énumération exacte des cinq lignes du catalogue qui satisfont \(\Psi(U)=010211\). Quatre lignes possèdent le bloc terminal \((169,272,272)\) ; la cinquième, qui porte le retour muté, possède \((418,674,521)\). La multiplicité maximale est quatre et n'est atteinte qu'une seule fois.
\end{proof}

La différence orientée entre le retour muté et le retour persistant est
\[
 \omega_\pi=b_\pi^- - b_\pi=(249,402,249),
 \qquad
 \langle(1,1,1),\omega_\pi\rangle=900.
\]
Cette identité sépare deux objets qu'il ne faut pas confondre. L'entier
\(169\) est une coordonnée d'une donnée terminale, c'est-à-dire une donnée de degré
zéro. Le vecteur \(\omega_\pi\) peut être lu comme une cochaîne de degré un
après orientation de l'arête qui relie les deux blocs. L'appeler cocycle exigerait,
en outre, de fixer un complexe de cochaînes et de vérifier que son cobord s'annule.
La prolongation analytique construite plus loin ne présuppose pas cette
identification.

La sélection de \(169\) n'exige pas de privilégier la première position du bloc.
Le groupe symétrique \(\mathfrak S_3\) agit en permutant ses trois coordonnées.
Le stabilisateur de \(b_\pi\) échange les deux occurrences de \(272\) et
fixe l'unique coordonnée de multiplicité un. Notons
\(\operatorname{sing}(b)\) cette coordonnée lorsque le bloc est de type
\((r,s,s)\), avec \(r\ne s\). De manière équivalente,
\[
\operatorname{sing}(b)
=\sum_{j=1}^3b_j-2\operatorname{moda}(b).
\]
En particulier,
\[
\boxed{\operatorname{sing}(b_\pi)=169.}
\]

\begin{definition}[Sélecteur résiduel de la microfibre]
\label{p3-83-c30-s1:def:selector-residual-pi}
Le sélecteur \(\rho_\pi\) applique successivement deux opérations :
\begin{enumerate}
  \item sélectionne le bloc terminal de multiplicité maximale au sein de \(p_{\rm ret}(\mathcal F_\pi)\) ;
  \item retient la coordonnée fixe sous l'action du stabilisateur de ce bloc.
\end{enumerate}
Par la \cref{p3-83-c30-s1:prop:bloque-terminal-persistente},
\[
\boxed{\rho_\pi=169.}
\]
\end{definition}

\begin{theorem}[Unicité équivariante du retour]
\label{p3-83-c30-s1:thm:unicidad-retorno-169}
Parmi les sélecteurs qui
\begin{enumerate}
  \item sont invariants sous les permutations des cinq régions ;
  \item sont équivariants sous les permutations des trois coordonnées
  terminales ;
  \item sélectionnent le bloc de multiplicité maximale ;
  \item retiennent une coordonnée fixe sous l'action de son stabilisateur,
\end{enumerate}
la valeur de retour est déterminée de manière unique et vaut \(169\).
\end{theorem}

\begin{proof}
L'invariance sous \(\mathfrak S_5\) empêche d'utiliser l'ordre des lignes. La
troisième condition sélectionne \(b_\pi\), qui est l'unique mode d'après la
\cref{p3-83-c30-s1:prop:bloque-terminal-persistente}. Son stabilisateur dans
\(\mathfrak S_3\) possède deux orbites sur les coordonnées : la paire de
valeurs \(272\) et la valeur singleton \(169\). La quatrième condition sélectionne la
seconde orbite. L'équivariance conserve sa valeur lorsque sa position est déplacée.
\end{proof}

Cet argument améliore la lecture positionnelle du retour. Le \(169\) n'est pas
« la première coordonnée » d'une ligne choisie : c'est un invariant du
multiensemble régional et de l'action de ses symétries de réordonnancement.
Il n'est pas non plus sélectionné pour sa rareté globale. Dans les \(468\) émissions du
catalogue, \(169\) apparaît \(84\) fois, uniformément réparti avec
\(14\) occurrences dans chacune des six positions. L'unicité du
\cref{p3-83-c30-s1:thm:unicidad-retorno-169} est donc relationnelle et régionale : elle réside
dans la combinaison de la microfibre, de la persistance modale et du stabilisateur, non dans la
fréquence isolée de l'entier.

% HMT_SOURCE_SEGMENT_END id=B02
\subsection{Construction du lecteur analytique régional et de son domaine de convergence}
\label{sec:s102:c38:2:3}

% HMT_SOURCE_SEGMENT_BEGIN id=B05 source=colaboracion/parte_iii/source/public_final/base_83/c30_body.tex body_lines=261-301

L'extension analytique est fixée par les règles suivantes.

\begin{definition}[Lecteur régional déterminantal]
\label{p3-83-c30-s1:def:lector-regional-determinantal}
Le lecteur régional déterminantal satisfait :
\begin{enumerate}
  \item \emph{compatibilité de degré} : la longueur cylindrique est le degré
  analytique;
  \item \emph{décomposition de renouvellement} : l'impulsion finie de retour et
  la continuation stricte appartiennent à des termes additifs distincts ;
  \item \emph{normalisation de continuation} : après l'enregistrement, une seule
  fois, du résidu \(\rho_\pi\), la branche survivante commence à l'unité de
  la droite déterminant ;
  \item \emph{covariance par concaténation} : chaque prolongation élémentaire
  agit au moyen de
  \(\bigwedge^2\mathcal T=D_A\) ;
  \item \emph{orientation} : dans la convention cardinale fixée pour APP, le retour régional appartient au secteur compensateur et est soustrait de la phase d'action de cinquième degré.
\end{enumerate}
\end{definition}

La troisième règle a un contenu mathématique. La valeur \(169\) est enregistrée
une fois comme amplitude du retour ; elle n'est pas multipliée à nouveau à chaque
prolongation. La continuation stricte compte l'unique droite qui se poursuit et
la normalise par son unité. C'est la différence entre une récurrence de
renouvellement et l'itération homogène de l'impulsion initiale.

La nécessité de déclarer cette normalisation peut être prouvée. Pour tout
\(\lambda\in\mathbb R\), la famille
\[
F_\lambda(z)
=169z^6+\lambda\frac{D_Az^7}{1-D_Az}
\]
conserve le même retour fini et le même rapport déterminantal entre
coefficients consécutifs. Le catalogue et le rapport seuls ne distinguent pas
\(\lambda\). L'unité de la droite déterminant fixe
\(\lambda=1\) avant d'évaluer \(z=\alpha\). De cette manière, la normalisation
ne s'obtient pas en ajustant la valeur finale ; elle fait partie de la définition du
caractère.

% HMT_SOURCE_SEGMENT_END id=B05
\subsection{Résolution analytique de la récurrence de retour}
\label{sec:s102:c38:2:4}

% HMT_SOURCE_SEGMENT_BEGIN id=B06 source=colaboracion/parte_iii/source/public_final/base_83/c30_body.tex body_lines=303-364

Nous définissons les coefficients \(\kappa_n\) par
\[
\kappa_6=169,
\qquad
\kappa_7=D_A,
\qquad
\kappa_{n+1}=D_A\kappa_n\quad(n\ge7).
\]
La première égalité est l'impulsion régionale ; les deux suivantes décrivent la
continuation normalisée.

\begin{theorem}[Réalisation analytique graduée]
\label{p3-83-c30-s1:thm:realizacion-analitica-graduada}
Le germe analytique compatible avec les règles de la
\cref{p3-83-c30-s1:def:lector-regional-determinantal} est unique et vaut
\[
\boxed{
\mathscr C_\pi(z)
=169z^6+\frac{D_Az^7}{1-D_Az}.}
\]
Ses coefficients satisfont
\[
\kappa_n=D_A^{n-6}\qquad(n\ge7).
\]
\end{theorem}

\begin{proof}
Écrivons
\[
F(z)=169z^6+\sum_{k\ge1}c_kz^{6+k}.
\]
La normalisation de continuation et la covariance déterminantale donnent
\[
c_1=D_A,
\qquad
c_{k+1}=D_Ac_k.
\]
Par récurrence, \(c_k=D_A^k\). La somme géométrique produit
\[
\sum_{k\ge1}D_A^kz^{6+k}
=\frac{D_Az^7}{1-D_Az}.
\]
Aucun coefficient libre ne subsiste.
\end{proof}

Comme \(0<\alpha<1\) et
\(0<D_A<1\),
\[
0<D_A\alpha<1.
\]
En conséquence, \(\mathscr C_\pi(\alpha)\) converge absolument et son
rayon de convergence est \(D_A^{-1}>1\). Numériquement,
\[
D_A=0.7751291192471564301888840964802\ldots,
\]
\[
D_A\alpha=0.00565639046986492673885079095469\ldots.
\]
La rapidité de cette contraction explique pourquoi la somme complète et sa
troncature de degré sept possèdent les mêmes quinze premiers chiffres dans la
carte finale.

% HMT_SOURCE_SEGMENT_END id=B06
\subsection{Persistance du retour dans les \texorpdfstring{\(5\)}{5} régions et compatibilité de la composante \texorpdfstring{\(\eta_{\mathrm{ret}}\)}{eta_ret}}
\label{sec:s102:c38:2:5}

% HMT_SOURCE_SEGMENT_BEGIN id=A02 source=colaboracion/parte_iii/source/public_final/ascensos_20260826/16f_realizacion_analitica_contraangulo.tex body_lines=49-170

Soit
\[
\mathcal F_\pi=\Psi^{-1}(010211)\subset\mathcal U_6
\]
la microfibre relative au catalogue. Rappelons ses cinq éléments, maintenant séparés en bloc initial et bloc terminal :
\[
\begin{aligned}
 &(501,614,498\mid169,272,272),\\
 &(810,923,870\mid169,272,272),\\
 &(810,983,810\mid169,272,272),\\
 &(870,923,810\mid169,272,272),\\
 &(870,923,810\mid418,674,521).
\end{aligned}
\]
Nous définissons la projection de retour
\[
p_{\rm ret}:\mathbb Z^6\longrightarrow\mathbb Z^3,
\qquad
p_{\rm ret}(u_1,\ldots,u_6)=(u_4,u_5,u_6).
\]

\begin{proposition}[Bloc terminal persistant]
\label{p3-20260826:prop:bloque-terminal-persistente}
Le multiensemble \(p_{\rm ret}(\mathcal F_\pi)\) possède un unique élément de
multiplicité maximale :
\[
b_\pi=(169,272,272),
\qquad
\operatorname{mult}_{\mathcal F_\pi}(b_\pi)=4.
\]
L'unique bloc terminal restant est
\[
b_\pi^-=(418,674,521)
\]
et a une multiplicité de un.
\end{proposition}

\begin{proof}
L'affirmation s'obtient par énumération exacte des cinq lignes du catalogue qui satisfont \(\Psi(U)=010211\). Quatre lignes possèdent le bloc terminal \((169,272,272)\) ; la cinquième, qui porte le retour muté, possède \((418,674,521)\). La multiplicité maximale est quatre et n'est atteinte qu'une seule fois.
\end{proof}

La différence orientée entre le retour muté et le retour persistant est
\[
 \omega_\pi=b_\pi^- - b_\pi=(249,402,249),
 \qquad
 \langle(1,1,1),\omega_\pi\rangle=900.
\]
Cette identité sépare deux objets qu'il ne faut pas confondre. L'entier
\(169\) est une coordonnée d'une donnée terminale, c'est-à-dire une donnée de degré
zéro. Le vecteur \(\omega_\pi\) peut être lu comme une cochaîne de degré un
après orientation de l'arête qui relie les deux blocs. L'appeler cocycle exigerait,
en outre, de fixer un complexe de cochaînes et de vérifier que son cobord s'annule.
La prolongation analytique construite plus loin ne présuppose pas cette
identification.

La sélection de \(169\) n'exige pas de privilégier la première position du bloc.
Le groupe symétrique \(\mathfrak S_3\) agit en permutant ses trois coordonnées.
Le stabilisateur de \(b_\pi\) échange les deux occurrences de \(272\) et
fixe l'unique coordonnée de multiplicité un. Notons
\(\operatorname{sing}(b)\) cette coordonnée lorsque le bloc est de type
\((r,s,s)\), avec \(r\ne s\). De manière équivalente,
\[
\operatorname{sing}(b)
=\sum_{j=1}^3b_j-2\operatorname{moda}(b).
\]
En particulier,
\[
\boxed{\operatorname{sing}(b_\pi)=169.}
\]

\begin{definition}[Sélecteur résiduel de la microfibre]
\label{p3-20260826:def:selector-residual-pi}
Le sélecteur \(\rho_\pi\) applique successivement deux opérations :
\begin{enumerate}
  \item sélectionne le bloc terminal de multiplicité maximale au sein de \(p_{\rm ret}(\mathcal F_\pi)\) ;
  \item retient la coordonnée fixe sous l'action du stabilisateur de ce bloc.
\end{enumerate}
Par la \cref{p3-20260826:prop:bloque-terminal-persistente},
\[
\boxed{\rho_\pi=169.}
\]
\end{definition}

\begin{theorem}[Unicité équivariante du retour]
\label{p3-20260826:thm:unicidad-retorno-169}
Parmi les sélecteurs qui
\begin{enumerate}
  \item sont invariants sous les permutations des cinq régions ;
  \item sont équivariants sous les permutations des trois coordonnées
  terminales ;
  \item sélectionnent le bloc de multiplicité maximale ;
  \item retiennent une coordonnée fixe sous l'action de son stabilisateur,
\end{enumerate}
la valeur de retour est déterminée de manière unique et vaut \(169\).
\end{theorem}

\begin{proof}
L'invariance sous \(\mathfrak S_5\) empêche d'utiliser l'ordre des lignes. La
troisième condition sélectionne \(b_\pi\), qui est l'unique mode d'après la
\cref{p3-20260826:prop:bloque-terminal-persistente}. Son stabilisateur dans
\(\mathfrak S_3\) possède deux orbites sur les coordonnées : la paire de
valeurs \(272\) et la valeur singleton \(169\). La quatrième condition sélectionne la
seconde orbite. L'équivariance conserve sa valeur lorsque sa position est déplacée.
\end{proof}

Cet argument améliore la lecture positionnelle du retour. Le \(169\) n'est pas
« la première coordonnée » d'une ligne choisie : c'est un invariant du
multiensemble régional et de l'action de ses symétries de réordonnancement.
Il n'est pas non plus sélectionné pour sa rareté globale. Dans les \(468\) émissions du
catalogue, \(169\) apparaît \(84\) fois, uniformément réparti avec
\(14\) occurrences dans chacune des six positions. L'unicité du
\cref{p3-20260826:thm:unicidad-retorno-169} est donc relationnelle et régionale : elle réside
dans la combinaison de la microfibre, de la persistance modale et du stabilisateur, non dans la
fréquence isolée de l'entier.

% HMT_SOURCE_SEGMENT_END id=A02
\subsection{Application du lecteur régional au couple angulaire conjugué}
\label{sec:s102:c38:2:6}

% HMT_SOURCE_SEGMENT_BEGIN id=A05 source=colaboracion/parte_iii/source/public_final/ascensos_20260826/16f_realizacion_analitica_contraangulo.tex body_lines=277-317

L'extension analytique est fixée par les règles suivantes.

\begin{definition}[Lecteur régional déterminantal]
\label{p3-20260826:def:lector-regional-determinantal}
Le lecteur régional déterminantal satisfait :
\begin{enumerate}
  \item \emph{compatibilité de degré} : la longueur cylindrique est le degré
  analytique;
  \item \emph{décomposition de renouvellement} : l'impulsion finie de retour et
  la continuation stricte appartiennent à des termes additifs distincts ;
  \item \emph{normalisation de continuation} : après l'enregistrement, une seule
  fois, du résidu \(\rho_\pi\), la branche survivante commence à l'unité de
  la droite déterminant ;
  \item \emph{covariance par concaténation} : chaque prolongation élémentaire
  agit au moyen de
  \(\bigwedge^2\mathcal T=D_A\) ;
  \item \emph{orientation} : dans la convention cardinale fixée pour APP, le retour régional appartient au secteur compensateur et est soustrait de la phase d'action de cinquième degré.
\end{enumerate}
\end{definition}

La troisième règle a un contenu mathématique. La valeur \(169\) est enregistrée
une fois comme amplitude du retour ; elle n'est pas multipliée à nouveau à chaque
prolongation. La continuation stricte compte l'unique droite qui se poursuit et
la normalise par son unité. C'est la différence entre une récurrence de
renouvellement et l'itération homogène de l'impulsion initiale.

La nécessité de déclarer cette normalisation peut être prouvée. Pour tout
\(\lambda\in\mathbb R\), la famille
\[
F_\lambda(z)
=169z^6+\lambda\frac{D_Az^7}{1-D_Az}
\]
conserve le même retour fini et le même rapport déterminantal entre
coefficients consécutifs. Le catalogue et le rapport seuls ne distinguent pas
\(\lambda\). L'unité de la droite déterminant fixe
\(\lambda=1\) avant d'évaluer \(z=\alpha\). De cette manière, la normalisation
ne s'obtient pas en ajustant la valeur finale ; elle fait partie de la définition du
caractère.

% HMT_SOURCE_SEGMENT_END id=A05
\subsection{Somme fermée et publication de la composante postérieure au retour}
\label{sec:s102:c38:2:7}

% HMT_SOURCE_SEGMENT_BEGIN id=A06 source=colaboracion/parte_iii/source/public_final/ascensos_20260826/16f_realizacion_analitica_contraangulo.tex body_lines=319-380

Nous définissons les coefficients \(\kappa_n\) par
\[
\kappa_6=169,
\qquad
\kappa_7=D_A,
\qquad
\kappa_{n+1}=D_A\kappa_n\quad(n\ge7).
\]
La première égalité est l'impulsion régionale ; les deux suivantes décrivent la
continuation normalisée.

\begin{theorem}[Réalisation analytique graduée]
\label{p3-20260826:thm:realizacion-analitica-graduada}
Le germe analytique compatible avec les règles de la
\cref{p3-20260826:def:lector-regional-determinantal} est unique et vaut
\[
\boxed{
\mathscr C_\pi(z)
=169z^6+\frac{D_Az^7}{1-D_Az}.}
\]
Ses coefficients satisfont
\[
\kappa_n=D_A^{n-6}\qquad(n\ge7).
\]
\end{theorem}

\begin{proof}
Écrivons
\[
F(z)=169z^6+\sum_{k\ge1}c_kz^{6+k}.
\]
La normalisation de continuation et la covariance déterminantale donnent
\[
c_1=D_A,
\qquad
c_{k+1}=D_Ac_k.
\]
Par récurrence, \(c_k=D_A^k\). La somme géométrique produit
\[
\sum_{k\ge1}D_A^kz^{6+k}
=\frac{D_Az^7}{1-D_Az}.
\]
Aucun coefficient libre ne subsiste.
\end{proof}

Comme \(0<\alpha<1\) et
\(0<D_A<1\),
\[
0<D_A\alpha<1.
\]
En conséquence, \(\mathscr C_\pi(\alpha)\) converge absolument et son
rayon de convergence est \(D_A^{-1}>1\). Numériquement,
\[
D_A=0.7751291192471564301888840964802\ldots,
\]
\[
D_A\alpha=0.00565639046986492673885079095469\ldots.
\]
La rapidité de cette contraction explique pourquoi la somme complète et sa
troncature de degré sept possèdent les mêmes quinze premiers chiffres dans la
carte finale.

% HMT_SOURCE_SEGMENT_END id=A06
\section{Confluence de \texorpdfstring{\(\eta_{\mathrm{ret}}\)}{eta_ret} avec la valuation \texorpdfstring{\(\varepsilon_H=-34\)}{epsilon_H=-34} dans la grandeur \texorpdfstring{\(\hbar_{\mathrm{ret}}\)}{hbar_ret}}
\label{sec:s102:c38:3}

\subsection{Dérivation du caractère déterminantal du transport parangulaire}
\label{sec:s102:c38:3:1}

% HMT_SOURCE_SEGMENT_BEGIN id=B04 source=colaboracion/parte_iii/source/public_final/base_83/c30_body.tex body_lines=196-259

La fermeture dodécaphasique fournit \(\alpha\). La complétion des trois
couples nonadiques fournit l'échelle \(10^3\), de sorte que
\[
A=10^3\alpha.
\]
La carte angulaire archimédienne étant choisie, la coordonnée commune en radians est
\[
x=\frac{\pi A}{180}=\frac{50\pi}{9}\alpha.
\]
Soit \(R\) une involution réelle de trace nulle en dimension deux,
\[
R^2=I_2,
\qquad
\operatorname{tr}R=0,
\]
et considérons le transport formel bicouche
\[
\mathcal T(x,y)=\exp(-xI_2-yR).
\]
La variable orientée \(y\) n'a pas encore été évaluée. Cependant, l'action
de \(\mathcal T\) sur la droite déterminante en est indépendante :
\[
\begin{aligned}
\det\mathcal T(x,y)
&=\exp\!\left(\operatorname{tr}(-xI_2-yR)\right)\\
&=e^{-2x}.
\end{aligned}
\]
On définit
\[
\boxed{
D_A=e^{-2x}=e^{-100\pi\alpha/9}.}
\]
L'indice rappelle qu'il s'agit du déterminant de la contraction commune associée à \(A\), et non d'une fonction de \(C_*^{\rm HMT}\).

\begin{proposition}[Caractère de la droite déterminante]
\label{p3-83-c30-s1:prop:caracter-determinante}
L'application
\[
\delta_A:(\mathbb N_0,+)\longrightarrow(\mathbb R_{>0},\cdot),
\qquad
\delta_A(k)=D_A^k,
\]
est l'unique caractère multiplicatif dont la valeur sur une prolongation
élémentaire est \(D_A\).
\end{proposition}

\begin{proof}
La multiplicativité donne
\[
\delta_A(k)
=\delta_A(\underbrace{1+\cdots+1}_{k\text{ fois}})
=\delta_A(1)^k=D_A^k.
\]
La même identité prouve l'existence et l'unicité.
\end{proof}

Le déterminant reste invariant par conjugaison de \(\mathcal T\) et par échange des feuillets \(R\mapsto-R\). La queue construite à partir de \(D_A\) appartient donc au mode commun. L'orientation apparaîtra dans le signe avec lequel le caractère régional corrige la phase.

% HMT_SOURCE_SEGMENT_END id=B04
\subsection{Caractère d'action de degré \texorpdfstring{\(5\)}{5} dans la construction régionale}
\label{sec:s102:c38:3:2}

% HMT_SOURCE_SEGMENT_BEGIN id=B07 source=colaboracion/parte_iii/source/public_final/base_83/c30_body.tex body_lines=366-405

Le bloc central
\[
\mathcal B_c=\begin{pmatrix}7&2\\2&7\end{pmatrix}
\]
possède les valeurs propres \(9\) et \(5\). Le cycle dodécaphasique apporte le rang
\(12\) ; ses orbites de pas trois possèdent une longueur de quatre ; et le recensement complet
apporte le quotient exact
\[
\frac{104976}{1944}=54.
\]
À partir de ces invariants, on fixe le caractère d'action de cinquième degré
\[
\boxed{
H_5(\alpha,\varphi)
=\frac12\sqrt{\frac{1000\alpha}{\varphi}}-\alpha
+\frac9{16}\alpha^2-\frac59\alpha^3
+\frac7{48}\alpha^4-\frac1{54}\alpha^5.}
\]
Les quatre coefficients rationnels peuvent s'écrire
\[
\frac9{16}=\frac{\lambda_+}{4^2},
\qquad
\frac59=\frac{\lambda_-}{\lambda_+},
\qquad
\frac7{48}=\frac{\tfrac12\operatorname{tr}\mathcal B_c}{4\cdot12},
\qquad
\frac1{54}=\frac{1944}{104976},
\]
avec \((\lambda_+,\lambda_-)=(9,5)\). Ces identités certifient la
provenance des nombres utilisés. L'affectation successive de ces
invariants aux degrés deux, trois, quatre et cinq fait partie du caractère
analytique ; elle ne se déduit pas de la seule existence du catalogue fini.

Cette précision élimine une ambiguïté habituelle. Une combinaison
d'invariants de la construction est une définition interne et reproductible ; sa
naturalité exige les règles du lecteur qui spécifient l'ordre, les signes et la
normalisation. Ici, ces règles sont visibles et sont soumises aux mêmes
contrôles négatifs que le reste de la construction.

% HMT_SOURCE_SEGMENT_END id=B07
\subsection{Décomposition exacte de la droite d'action en sections pré-retour et post-retour}
\label{sec:s102:c38:3:3}

% HMT_SOURCE_SEGMENT_BEGIN id=B12 source=colaboracion/parte_iii/source/public_final/base_83/c30_body.tex body_lines=571-617

Le caractère de cinquième degré et la composante d'action postérieure au retour
sont deux valeurs distinctes :
\[
H_5
=1.05457181764615446962582182943306\ldots,
\]
\[
\bar h_{\rm HMT}^{\rm ret}
=1.05457181691503906113369058472430\ldots.
\]
La première est l'action locale avant l'incorporation du caractère régional ; la
seconde est l'action orientée qui reste après le retour. La correction
qui engendre la coordonnée angulaire conjuguée empêche de les identifier.

La comparaison métrologique doit respecter cette distinction. La première valeur donne
la mantisse
\[
2\pi H_5
=6.62607014999998792600111034989947\ldots,
\]
tandis que l'action postérieure donne
\[
2\pi\bar h_{\rm HMT}^{\rm ret}
=6.62607014540625433351074984759288\ldots.
\]
Dans le Système international, la valeur
\[
h=6.62607015\times10^{-34}\;\mathrm{J\,s}
\]
est exacte par définition \cite{BIPMSI}. C'est pourquoi,
\[
2\pi H_5-6.62607015
=-1.2073998889650101\ldots\times10^{-14},
\]
et ce n'est pas une égalité exacte. La proximité de la valeur antérieure au retour est
un contrôle numérique remarquable ; elle n'autorise pas à remplacer la constante définitoire
du SI par la composante postérieure ni à confondre les deux actions HMT\@.

Ce résultat détermine avec précision la situation mathématique. Le lecteur
régional engendre \(C_*^{\rm HMT}\) et, à partir de celui-ci, la fonctionnelle
\(\Gamma_{6\to8}\). Le même calcul sépare l'approximation de Planck
produite par \(H_5\) de l'action réduite qui intervient dans
la coordonnée angulaire conjuguée. La séparation n'affaiblit pas la construction : elle élimine une
identification incompatible avec les chiffres exacts et transforme chaque
comparaison en une affirmation réfutable de type bien défini.

% HMT_SOURCE_SEGMENT_END id=B12
\subsection{Factorisation du caractère déterminantal dans les \texorpdfstring{\(2\)}{2} cartes de phase}
\label{sec:s102:c38:3:4}

% HMT_SOURCE_SEGMENT_BEGIN id=A04 source=colaboracion/parte_iii/source/public_final/ascensos_20260826/16f_realizacion_analitica_contraangulo.tex body_lines=211-275

La clôture dodécaphasique fournit \(\alpha\). L'insertion angulaire déclarée
emploie l'échelle \(10^3\), compatible avec le regroupement de trois paires
nonadiques, de sorte que
\[
\Adeg=10^3\alpha\,{}^\circ.
\]
La carte angulaire archimédienne étant choisie, la coordonnée commune en radians est
\[
x=: \Arad=\frac{\pi\Adeg}{180}=\frac{50\pi}{9}\alpha.
\]
Soit \(R\) une involution réelle de trace nulle en dimension deux,
\[
R^2=I_2,
\qquad
\operatorname{tr}R=0,
\]
et considérons le transport formel bicouche
\[
\mathcal T(x,y)=\exp(-xI_2-yR).
\]
La variable orientée \(y\) n'a pas encore été évaluée. Cependant, l'action
de \(\mathcal T\) sur la droite déterminante en est indépendante :
\[
\begin{aligned}
\det\mathcal T(x,y)
&=\exp\!\left(\operatorname{tr}(-xI_2-yR)\right)\\
&=e^{-2x}.
\end{aligned}
\]
On définit
\[
\boxed{
D_A=e^{-2x}=e^{-100\pi\alpha/9}.}
\]
L'indice rappelle qu'il s'agit du déterminant de la contraction commune
associée à \(A\), non d'une fonction du contre-angle.

\begin{proposition}[Caractère de la droite déterminante]
\label{p3-20260826:prop:caracter-determinante}
L'application
\[
\delta_A:(\mathbb N_0,+)\longrightarrow(\mathbb R_{>0},\cdot),
\qquad
\delta_A(k)=D_A^k,
\]
est l'unique caractère multiplicatif dont la valeur sur une prolongation
élémentaire est \(D_A\).
\end{proposition}

\begin{proof}
La multiplicativité donne
\[
\delta_A(k)
=\delta_A(\underbrace{1+\cdots+1}_{k\text{ fois}})
=\delta_A(1)^k=D_A^k.
\]
La même identité prouve l'existence et l'unicité.
\end{proof}

Le déterminant reste invariant par conjugaison de \(\mathcal T\) et par échange des feuillets \(R\mapsto-R\). La queue construite à partir de \(D_A\) appartient donc au mode commun. L'orientation apparaîtra dans le signe avec lequel le caractère régional corrige la phase.

% HMT_SOURCE_SEGMENT_END id=A04
\subsection{Publication du caractère d'action de degré \texorpdfstring{\(5\)}{5} dans la carte parangulaire}
\label{sec:s102:c38:3:5}

% HMT_SOURCE_SEGMENT_BEGIN id=A07 source=colaboracion/parte_iii/source/public_final/ascensos_20260826/16f_realizacion_analitica_contraangulo.tex body_lines=382-418

Le bloc central
\[
\mathcal B_c=\begin{pmatrix}7&2\\2&7\end{pmatrix}
\]
possède les valeurs propres \(9\) et \(5\). Le cycle dodécaphasique apporte le rang
\(12\) ; ses orbites de pas trois possèdent une longueur de quatre ; et le recensement complet
apporte le quotient exact
\[
\frac{104976}{1944}=54.
\]
Le lecteur en vigueur déclare le caractère d'action de cinquième degré
\[
\boxed{
H_5(\alpha,\varphi)
=\frac12\sqrt{\frac{1000\alpha}{\varphi}}-\alpha
+\frac9{16}\alpha^2-\frac59\alpha^3
+\frac7{48}\alpha^4-\frac1{54}\alpha^5.}
\]
Les quatre coefficients rationnels peuvent s'écrire
\[
\frac9{16}=\frac{\lambda_+}{4^2},
\qquad
\frac59=\frac{\lambda_-}{\lambda_+},
\qquad
\frac7{48}=\frac{\tfrac12\operatorname{tr}\mathcal B_c}{4\cdot12},
\qquad
\frac1{54}=\frac{1944}{104976},
\]
avec \((\lambda_+,\lambda_-)=(9,5)\). Ces identités fixent la
représentabilité arithmétique interne des coefficients du caractère. Leur
ordre, leurs signes et leur normalisation appartiennent à la règle graduée du lecteur
régional et sont appliqués avant toute carte d'unités. Les ablations
certifient que chaque terme intervient effectivement dans la phase résultante.
En particulier, ni la mantisse SI de l'action ni la valeur sexagésimale du
contre-angle ne font partie du domaine de la fonction génératrice.

% HMT_SOURCE_SEGMENT_END id=A07
\section{Application parangulaire \texorpdfstring{\(P(\alpha,\eta_{\mathrm{ret}})\)}{P(alpha,eta_ret)} et clôture bicouche \texorpdfstring{\(C^*=2(\eta_{\mathrm{ret}}+\alpha)\)}{C*=2(eta_ret+alpha)}}
\label{sec:s102:c38:4}

\subsection{Phase orientée et coordonnée angulaire conjuguée}
\label{sec:s102:c38:4:1}

% HMT_SOURCE_SEGMENT_BEGIN id=B08 source=colaboracion/parte_iii/source/public_final/base_83/c30_body.tex body_lines=407-472

La phase de cinquième degré antérieure au retour est
\[
y_5=\frac\pi{90}\bigl(H_5+\alpha\bigr).
\]
Le caractère régional complet produit la correction
\(\mathscr C_\pi(\alpha)\). Avec l'orientation fixée dans la
\cref{p3-83-c30-s1:def:lector-regional-determinantal}, nous définissons
\[
\boxed{
y_*=\frac\pi{90}\bigl(H_5+\alpha\bigr)
-169\alpha^6
-\frac{D_A\alpha^7}{1-D_A\alpha}.}
\]
C'est la formule primaire. L'unité sexagésimale n'intervient qu'à l'application de la carte
\[
C_*^{\rm HMT}=\frac{180}{\pi}y_*.
\]

\begin{theorem}[Coordonnée angulaire conjuguée du lecteur régional déterminantal]
\label{p3-83-c30-s1:thm:contraangulo-regional}
Une fois fixés la clôture dodécaphasique de \(\alpha\),
\(\varphi=(1+\sqrt5)/2\) et les règles du lecteur régional déterminantal, la
phase \(y_*\) et sa coordonnée angulaire conjuguée sont déterminées sans introduire la valeur
recherchée. Leur évaluation est
\[
\boxed{
y_*=0.03706622646583826398493779409230638\ldots,}
\]
\[
\boxed{
C_*^{\rm HMT}
=2.12373833896864572426195138039336\ldots{}^\circ.}
\]
Donc, à quinze décimales,
\[
\boxed{C_*^{\rm HMT}=2.123738338968646^\circ.}
\]
\end{theorem}

\begin{proof}
Le \cref{p3-83-c30-s1:thm:unicidad-retorno-169} détermine \(169\). La clôture de
\(\alpha\) détermine \(A\), \(x\) et
\(D_A=e^{-100\pi\alpha/9}\). Le \cref{p3-83-c30-s1:thm:realizacion-analitica-graduada}
détermine la série complète. La formule de \(H_5\) contient uniquement
\(\alpha\), \(\varphi\) et des invariants structurels déclarés. Substituer
ces données dans l'expression de \(y_*\) donne la valeur imprimée. La valeur décimale
de \(C_*^{\rm HMT}\) n'apparaît dans aucun des membres de droite.
\end{proof}

La composante réduite d'action postérieure au retour est désormais une
sortie :
\[
\boxed{
\bar h_{\rm HMT}^{\rm ret}
=\frac{C_*^{\rm HMT}}2-\alpha
=H_5-\frac{90}{\pi}\mathscr C_\pi(\alpha).}
\]
Numériquement,
\[
\bar h_{\rm HMT}^{\rm ret}
=1.05457181691503906113369058472430\ldots.
\]
L'identité inverse n'a pas été utilisée pour introduire \(C_*^{\rm HMT}\) ; elle
s'obtient après la génération de \(y_*\).

% HMT_SOURCE_SEGMENT_END id=B08
\subsection{Coordonnée conjuguée \texorpdfstring{\(C^*\)}{C*} et réversion de son orientation sous l'involution bicouche}
\label{sec:s102:c38:4:2}

% HMT_SOURCE_SEGMENT_BEGIN id=A08 source=colaboracion/parte_iii/source/public_final/ascensos_20260826/16f_realizacion_analitica_contraangulo.tex body_lines=420-503

La phase de cinquième degré antérieure au retour est
\[
y_5=\frac\pi{90}\bigl(H_5+\alpha\bigr).
\]
Le caractère régional complet produit la correction
\(\mathscr C_\pi(\alpha)\). Avec l'orientation fixée dans la
\cref{p3-20260826:def:lector-regional-determinantal}, nous définissons
\[
\boxed{
\Crad:=y_*=\frac\pi{90}\bigl(H_5+\alpha\bigr)
-169\alpha^6
-\frac{D_A\alpha^7}{1-D_A\alpha}.}
\]
Désormais, nous écrivons aussi
\[
 \boxed{y_{\rm reg}:=y_*}
\]
lorsqu'il est nécessaire de distinguer cette réalisation de la phase spectrale de
\Gtr. Les deux notations désignent ici le même objet régional canonique.
Telle est la formule primaire. L'unité sexagésimale n'intervient qu'à l'application de la carte
\[
\boxed{\Cdeg:=\frac{180}{\pi}\Crad.}
\]
Les symboles \(\Crad\) et \(\Cdeg\) ne désignent pas deux contre-angles : ce sont les
coordonnées radiale et sexagésimale du même élément angulaire.

\begin{theorem}[Contre-angle du lecteur régional déterminantal]
\label{p3-20260826:thm:contraangulo-regional}
La clôture dodécaphasique de \(\alpha\),
\(\varphi=(1+\sqrt5)/2\), le germe \(H_5\) et les règles du lecteur régional
déterminantal étant fixés, la phase \(\Crad\) et son contre-angle sont déterminés sans
introduire la valeur recherchée dans la fonction d'évaluation. Son évaluation est
\[
\boxed{
\Crad=0.03706622646583826398493779409230638\ldots,}
\]
\[
\boxed{
\Cdeg
=2.12373833896864572426195138039336\ldots{}^\circ.}
\]
Donc, à quinze décimales,
\[
\boxed{\Cdeg=2.123738338968646^\circ.}
\]
\end{theorem}

\begin{proof}
Le \cref{p3-20260826:thm:unicidad-retorno-169} détermine \(169\). La clôture de
\(\alpha\) détermine \(\Adeg\), \(\Arad\) et
\(D_A=e^{-100\pi\alpha/9}\). Le \cref{p3-20260826:thm:realizacion-analitica-graduada}
détermine la série complète. La formule fixée de \(H_5\) contient uniquement
\(\alpha\), \(\varphi\) et des invariants structurels déclarés. Substituer
ces données dans l'expression de \(\Crad\) donne la valeur imprimée. La valeur décimale
du contre-angle n'apparaît dans aucun des membres de droite de cette
évaluation. La comparaison avec une mantisse métrologique d'action est effectuée
uniquement après l'obtention de la phase et ne modifie pas la fonction génératrice.
\end{proof}

Soit \(u_\circ\) la base sexagésimale dont le tour complet a pour coordonnée
\(360\), et introduisons la coordonnée angulaire
\(\alpha_{\angle,{\rm deg}}:=\alpha_{\rm HMT}\). La sortie postérieure au
retour est alors typée par
\[
\boxed{
\etaRet
=\frac{\Cdeg}2-\alpha_{\angle,{\rm deg}}
=H_5-\frac{90}{\pi}\mathscr C_\pi(\alpha).}
\]
Sa coordonnée radiale est
\[
 \etaRetrad=\frac{\pi}{180}\etaRet
 =\frac{\Crad}{2}-\frac{\pi}{180}\alpha_{\rm HMT}.
\]
Numériquement,
\[
\etaRet
=1.05457181691503906113369058472430\ldots.
\]
L'identité inverse n'a pas été utilisée pour introduire \(\Cdeg\) ; elle s'obtient
après la génération de \(\Crad\). Le symbole \(\hbar\) reste réservé à
l'élément de la droite d'action construit dans la section suivante.

% HMT_SOURCE_SEGMENT_END id=A08
\subsection{Théorème direct \texorpdfstring{\(C^*=2(\eta_{\mathrm{ret}}+\alpha)\)}{C*=2(eta_ret+alpha)} et caractère inverse de l'identité \texorpdfstring{\(\eta_{\mathrm{ret}}=C^*/2-\alpha\)}{eta_ret=C*/2-alpha}}
\label{sec:s102:c38:4:3}

% HMT_SOURCE_SEGMENT_BEGIN id=A13 source=colaboracion/parte_iii/source/public_final/ascensos_20260826/16f_realizacion_analitica_contraangulo.tex body_lines=750-797

Le caractère de cinquième degré et la coordonnée postérieure au retour
sont deux valeurs distinctes :
\[
H_5
=1.05457181764615446962582182943306\ldots,
\]
\[
\etaRet
=1.05457181691503906113369058472430\ldots.
\]
Le premier est le caractère local avant l'intégration du caractère régional ; le
second est la coordonnée orientée qui subsiste après le retour. La correction
qui génère le contre-angle empêche de les identifier.

La comparaison métrologique doit respecter cette distinction. La première valeur donne
la mantisse
\[
2\pi H_5
=6.62607014999998792600111034989947\ldots,
\]
tandis que la coordonnée ultérieure donne
\[
2\pi\etaRet
=6.62607014540625433351074984759288\ldots.
\]
À titre de contrôle externe uniquement, nous retenons la mantisse \(6.62607015\) de la
valeur définitoire du Système international \cite{BIPMSI}. On n'intègre
pas encore son exposant ni son unité à la droite HMT : la décennie est dérivée dans le
\cref{chap:orden-determinantal-accion} et l'unité appartient à une carte
métrologique ultérieure. La comparaison des mantisses donne
\[
2\pi H_5-6.62607015
=-1.2073998889650101\ldots\times10^{-14},
\]
et ce n'est pas une égalité exacte. La proximité de la valeur antérieure au retour est
un contrôle numérique remarquable ; elle n'autorise pas à remplacer la constante définitoire
du SI par la coordonnée ultérieure ni à confondre une coordonnée angulaire avec
l'élément d'action qui sera construit lors de l'intégration de la décennie.

Ce résultat détermine avec précision la situation mathématique. Le lecteur
régional génère \((\Crad,\Cdeg)\) et, à partir de cet unique élément angulaire, la fonctionnelle
\(\Gamma_{6\to8}\). Le même calcul sépare la proximité externe de la
mantisse produite par \(H_5\) de la coordonnée de retour qui intervient dans le
contre-angle. La séparation n'affaiblit pas la construction : elle élimine une
identification incompatible avec les chiffres exacts et transforme chaque
comparaison en une affirmation falsifiable de type bien défini.

% HMT_SOURCE_SEGMENT_END id=A13
% HMT_SOURCE_SEGMENT_BEGIN id=D01 source=manuscrito/sucesor_102/deltas_ley9/c38_par_angular_cartas_y_canales.tex body_lines=3-20
\label{sec:l9s102:par-angular}

La complétion cubique
\begin{equation}
 10^3=(9+1)^3=9^3+3\cdot9^2+3\cdot9+1
 =729+243+27+1
 \label{eq:l9s102:base-1000}
\end{equation}
détermine la carte globale de la coordonnée directe. La clôture bicouche apporte
le facteur \(2\) de la coordonnée conjuguée :
\begin{equation}
 \boxed{
 A_{\deg}=10^3\alpha_{\mathrm{HMT}},
 \qquad
 C^*_{\deg}=2(\eta_{\mathrm{ret}}+\alpha_{\mathrm{HMT}})
 \quad\text{dans }\mathbb R/(360\mathbb Z).}
 \label{eq:l9s102:a-cstar-deg}
\end{equation}
% HMT_SOURCE_SEGMENT_END id=D01
% HMT_SOURCE_SEGMENT_BEGIN id=D04 source=manuscrito/sucesor_102/deltas_ley9/c38_par_angular_cartas_y_canales.tex body_lines=60-66
L'identité inverse
\begin{equation}
 \eta_{\mathrm{ret}}=\frac12C^*_{\deg}-\alpha_{\mathrm{HMT}}
 \label{eq:l9s102:eta-corolario-inverso}
\end{equation}
est établie comme corollaire algébrique de la construction directe.

% HMT_SOURCE_SEGMENT_END id=D04
\section{Commutativité exacte des cartes en tours, degrés et radians}
\label{sec:s102:c38:5}

\subsection{Dérivation de la longueur cylindrique et du degré analytique du précurseur parangulaire}
\label{sec:s102:c38:5:1}

% HMT_SOURCE_SEGMENT_BEGIN id=B03 source=colaboracion/parte_iii/source/public_final/base_83/c30_body.tex body_lines=157-194

Soit \(\mathcal P\) le module libre engendré par les mots régionaux et
soit \(F^n\mathcal P\) sa filtration par longueur. La graduation associée
enregistre le premier niveau auquel apparaît chaque mot. Pour un paramètre
\(z\), le caractère ordinaire de longueur est
\[
\chi_z([w])=z^{|w|}.
\]
La multiplicativité
\[
\chi_z([uv])=\chi_z([u])\chi_z([v])
\]
exprime que la concaténation additionne les longueurs. Lors de l'évaluation en
\(z=\alpha\), un mot de longueur \(n\) reçoit le degré \(\alpha^n\).

\begin{definition}[Compatibilité de degré]
\label{p3-83-c30-s1:def:compatibilidad-grado}
Le lecteur analytique régional est compatible avec la filtration cylindrique
lorsqu'il envoie la composante homogène de longueur \(n\) sur le degré analytique
\(n\) :
\[
\operatorname{gr}_n\mathcal P\longrightarrow
\alpha^n\mathbb R.
\]
\end{definition}

Les cinq régions appartiennent à \(\mathcal U_6\) ; c'est pourquoi la donnée de
retour apparaît en degré six :
\[
\rho_\pi[U_6]\longmapsto169\alpha^6.
\]
Ce six est la longueur du mot régional. Ce n'est pas la longueur des
marches fermées du \cref{chap:cinco-ciclos-pi}. Ces marches parcourent à
la fois les quatre ancrages
\(U_\pi,U_e,U_\varphi,U_{\rm APP}\) et possèdent une longueur minimale de huit. Un
mot, une arête du graphe de blocs et une marche fermée sont des objets de
types distincts ; la graduation utilisée ici appartient au premier.

% HMT_SOURCE_SEGMENT_END id=B03
\subsection{Publication du degré analytique dans la carte globale de phase}
\label{sec:s102:c38:5:2}

% HMT_SOURCE_SEGMENT_BEGIN id=A03 source=colaboracion/parte_iii/source/public_final/ascensos_20260826/16f_realizacion_analitica_contraangulo.tex body_lines=172-209

Soit \(\mathcal P\) le module libre engendré par les mots régionaux et
soit \(F^n\mathcal P\) sa filtration par longueur. La graduation associée
enregistre le premier niveau auquel apparaît chaque mot. Pour un paramètre
\(z\), le caractère ordinaire de longueur est
\[
\chi_z([w])=z^{|w|}.
\]
La multiplicativité
\[
\chi_z([uv])=\chi_z([u])\chi_z([v])
\]
exprime que la concaténation additionne les longueurs. Lors de l'évaluation en
\(z=\alpha\), un mot de longueur \(n\) reçoit le degré \(\alpha^n\).

\begin{definition}[Compatibilité de degré]
\label{p3-20260826:def:compatibilidad-grado}
Le lecteur analytique régional est compatible avec la filtration cylindrique
lorsqu'il envoie la composante homogène de longueur \(n\) sur le degré analytique
\(n\) :
\[
\operatorname{gr}_n\mathcal P\longrightarrow
\alpha^n\mathbb R.
\]
\end{definition}

Les cinq régions appartiennent à \(\mathcal U_6\) ; c'est pourquoi la donnée de
retour apparaît en degré six :
\[
\rho_\pi[U_6]\longmapsto169\alpha^6.
\]
Ce six est la longueur du mot régional. Ce n'est pas la longueur des
marches fermées du \cref{chap:cinco-ciclos-pi}. Ces marches parcourent à
la fois les quatre ancrages
\(U_\pi,U_e,U_\varphi,U_{\rm APP}\) et possèdent une longueur minimale de huit. Un
mot, une arête du graphe de blocs et une marche fermée sont des objets de
types distincts ; la graduation utilisée ici appartient au premier.

% HMT_SOURCE_SEGMENT_END id=A03
% HMT_SOURCE_SEGMENT_BEGIN id=D02 source=manuscrito/sucesor_102/deltas_ley9/c38_par_angular_cartas_y_canales.tex body_lines=21-34
Le calendrier global présente les trois factorisations
\begin{equation}
 360=12\cdot30=4\cdot90=3\cdot120.
 \label{eq:l9s102:360-factorizaciones}
\end{equation}
La carte analytique s'obtient au moyen de
\begin{equation}
 A_{\mathrm{rad}}=\frac\pi{180}A_{\deg}
 =\frac{50\pi}{9}\alpha_{\mathrm{HMT}},
 \qquad
 C^*_{\mathrm{rad}}=\frac\pi{180}C^*_{\deg}
 =\frac\pi{90}(\eta_{\mathrm{ret}}+\alpha_{\mathrm{HMT}}).
 \label{eq:l9s102:a-cstar-rad}
\end{equation}
% HMT_SOURCE_SEGMENT_END id=D02
% HMT_SOURCE_SEGMENT_BEGIN id=D03 source=manuscrito/sucesor_102/deltas_ley9/c38_par_angular_cartas_y_canales.tex body_lines=35-59

\begin{theorem}[Commutativité des factorisations régionale et d'action]
\label{thm:l9s102:conmutatividad-cstar}
Les deux chemins vers la phase conjuguée produisent la même sortie :
\begin{equation}
 \boxed{
 \frac\pi{180}C^*_{\deg}
 =\frac\pi{90}(\eta_{\mathrm{ret}}+\alpha_{\mathrm{HMT}})
 =\frac\pi{90}(H_5+\alpha_{\mathrm{HMT}})
 -\mathscr C_\pi(\alpha_{\mathrm{HMT}})=y_*.}
 \label{eq:l9s102:cuadrado-cstar}
\end{equation}
\end{theorem}

\begin{proof}
En substituant \cref{eq:l9s102:eta-ret} dans le second membre, on obtient
\[
 \frac\pi{90}
 \left(H_5-\frac{90}{\pi}\mathscr C_\pi+\alpha_{\mathrm{HMT}}\right)
 =\frac\pi{90}(H_5+\alpha_{\mathrm{HMT}})-\mathscr C_\pi.
\]
La dernière expression est la définition régionale de \(y_*\), tandis que la
première égalité procède de \cref{eq:l9s102:a-cstar-deg}.
\end{proof}

% HMT_SOURCE_SEGMENT_END id=D03
\section{Décomposition spectrale \texorpdfstring{\(q_\pm\)}{q+/-} et calcul fonctionnel sur les canaux \texorpdfstring{\(90/120\)}{90/120} préalablement générés}
\label{sec:s102:c38:6}

\subsection{Reconstruction spectrale du couple angulaire et transition hexade--octade}
\label{sec:s102:c38:6:1}

% HMT_SOURCE_SEGMENT_BEGIN id=B11 source=colaboracion/parte_iii/source/public_final/base_83/c30_body.tex body_lines=523-569

La phase commune et la phase orientée déterminent les deux canaux
\[
q_-=e^{-x+y_*},
\qquad
q_+=e^{-x-y_*}.
\]
On retrouve exactement les invariants
\[
q_-q_+=D_A,
\qquad
\alpha=-\frac9{100\pi}\log D_A,
\qquad
C_*^{\rm HMT}=\frac{90}{\pi}\log\frac{q_-}{q_+}.
\]
La première égalité lit la contraction commune ; la troisième, la séparation
orientée des feuillets. L'échange \(R\mapsto-R\) permute
\(q_+\) et \(q_-\), conserve \(D_A\) et inverse le signe de \(y_*\).

Pour
\[
S_{90}(q)=\frac{q^{90}}{1-q^{270}},
\qquad
S_{120}(q)=\frac{q^{120}}{1-q^{360}},
\]
la fonctionnelle de transition hexade--octade s'écrit intégralement dans la
carte des phases :
\[
\Gamma_{6\to8}
=\frac{180}{\pi}xy_*
+12\bigl(S_{90}(q_-)-S_{90}(q_+)\bigr)
-\bigl(S_{120}(q_-)-S_{120}(q_+)\bigr).
\]
L'évaluation produite par la coordonnée \(C_*^{\rm HMT}\) du
\cref{p3-83-c30-s1:thm:contraangulo-regional} est
\[
\boxed{
\Gamma_{6\to8}
=0.274007672694459274747255260774\ldots.}
\]
Ainsi, la valeur interne reconnue en mécanique dimensionnelle comme fonctionnelle de
Barbero--Immirzi cesse de recevoir \(C_*^{\rm HMT}\) comme littéral indépendant :
toutes deux descendent du même couple de phases engendré dans ce chapitre.
L'interprétation géométrique et physique de \(\Gamma_{6\to8}\) appartient à l'œuvre
de mécanique dimensionnelle ; son évaluation mathématique appartient déjà à l'interface
précédemment construite.

% HMT_SOURCE_SEGMENT_END id=B11
\subsection{Transporteur relatif et disjonction spectrale des blocs}
\label{sec:s102:c38:6:2}

% HMT_SOURCE_SEGMENT_BEGIN id=B13 source=colaboracion/parte_iii/source/public_final/base_83/c30_body.tex body_lines=619-675
\label{p3-83-c30-s1:sec:transportador-relativo-split}

Dans la base ordonnée des canaux, soit
\[
 R=\begin{pmatrix}0&1\\1&0\end{pmatrix},
 \qquad P_\pm=\frac{I\pm R}{2},
 \qquad
 B_0=7I+2R,
 \qquad B_1=\Adeg I+\Cdeg R.
\]
La décomposition scindée
\[
 xI+yR=\rho(x,y)e^{\eta(x,y)R},
 \quad \rho(x,y)=\sqrt{x^2-y^2},
 \quad \eta(x,y)=\operatorname{artanh}(y/x)
\]
construit l'unique transporteur
\[
 T_{\rm rel}
 =\frac{\sqrt{(\Adeg)^2-(\Cdeg)^2}}{\sqrt{45}}
   \exp\!\left[
   \left(\operatorname{artanh}\frac{\Cdeg}{\Adeg}
   -\operatorname{artanh}\frac27\right)R\right],
 \qquad T_{\rm rel}B_0=B_1.
\]
Dans la base spectrale, il agit par
\(9\mapsto\Adeg+\Cdeg\) et
\(5\mapsto\Adeg-\Cdeg\).

\begin{theorem}[Disjonction spectrale des réalisations locale et angulaire]
\label{p3-83-c30-s1:thm:no-go-entrelazador-bloques-split}
Les spectres
\[
 \operatorname{Spec}(B_0)=\{9,5\},\qquad
 \operatorname{Spec}(B_1)=\{\Adeg+\Cdeg,\Adeg-\Cdeg\}
\]
sont disjoints. En conséquence,
\[
 \operatorname{Hom}_{B_0,B_1}
 :=\{\Phi:\Phi B_0=B_1\Phi\}=\{0\}.
\]
\end{theorem}

\begin{proof}
Les bornes rationnelles construites pour les coordonnées angulaires donnent
\(7.29<\Adeg<7.30\) et \(2.12<\Cdeg<2.13\) ; donc,
\(9.41<\Adeg+\Cdeg<9.43\) et
\(5.16<\Adeg-\Cdeg<5.18\). Dans la base qui diagonalise \(R\), chaque entrée
d'un entrelaceur est multipliée par la différence entre une valeur propre
de \(B_0\) et une de \(B_1\) ; la disjonction impose qu'elles soient toutes nulles.
\end{proof}

Le transporteur relatif compare deux blocs au sein de l'algèbre scindée ;
les entrelaceurs dodécaphasique--Witt et d'incidence agissent entre des
représentations équivalentes dans leurs espaces de résidence propres. Cette séparation
type positivement les deux opérations et conserve leurs domaines.

% HMT_SOURCE_SEGMENT_END id=B13
\subsection{Évaluation spectrale \texorpdfstring{\(V_{90,120}\)}{V_90,120} des canaux d'incidence préalablement générés par le TPK}
\label{sec:s102:c38:6:3}

% HMT_SOURCE_SEGMENT_BEGIN id=A11 source=colaboracion/parte_iii/source/public_final/ascensos_20260826/16f_realizacion_analitica_contraangulo.tex body_lines=554-572

La phase commune et la phase orientée déterminent les deux canaux
\[
q_-=e^{-\Arad+\Crad},
\qquad
q_+=e^{-\Arad-\Crad}.
\]
On retrouve exactement les invariants
\[
q_-q_+=D_A,
\qquad
\alpha=-\frac9{100\pi}\log D_A,
\qquad
\Cdeg=\frac{90}{\pi}\log\frac{q_-}{q_+}.
\]
La première égalité lit la contraction commune ; la troisième, la séparation
orientée des feuillets. L'échange \(R\mapsto-R\) permute
\(q_+\) et \(q_-\), conserve \(D_A\) et inverse le signe de \(\Crad\).

% HMT_SOURCE_SEGMENT_END id=A11
% HMT_SOURCE_SEGMENT_BEGIN id=D05 source=manuscrito/sucesor_102/deltas_ley9/c38_par_angular_cartas_y_canales.tex body_lines=68-97
\label{sec:l9s102:qpm-posteriores}

Soient \(H^2=I\) et \(P_\pm=(I\pm H)/2\). Le bloc angulaire
\begin{equation}
 B=A_{\mathrm{rad}}I+C^*_{\mathrm{rad}}H
 =(A_{\mathrm{rad}}+C^*_{\mathrm{rad}})P_+
 +(A_{\mathrm{rad}}-C^*_{\mathrm{rad}})P_-
 \label{eq:l9s102:b-angular}
\end{equation}
se publie multiplicativement par
\begin{equation}
 e^{-B}=q_+P_++q_-P_-,
 \qquad
 q_\pm=\exp[-(A_{\mathrm{rad}}\pm C^*_{\mathrm{rad}})].
 \label{eq:l9s102:qpm-angular}
\end{equation}
L'involution \(C^*\mapsto-C^*\) permute \(q_+\) et \(q_-\), et conserve
\begin{equation}
 q_+q_-=e^{-2A_{\mathrm{rad}}}=D_A.
 \label{eq:l9s102:det-angular}
\end{equation}
Les cardinaux \(90/120\), construits préalablement par l'incidence du
TPK, reçoivent maintenant leur évaluation spectrale :
\begin{equation}
 V_{90,120}
 =(I-T^{90})(I-T^{120})^{-1}
 =\sum_{\sigma\in\{+,-\}}
 \frac{1-q_\sigma^{90}}{1-q_\sigma^{120}}P_\sigma.
 \label{eq:l9s102:v90120-posterior}
\end{equation}
% HMT_SOURCE_SEGMENT_END id=D05
\subsubsection{Transporteur relatif scindé}
\label{sec:s102:c38:6:3:1}

% HMT_SOURCE_SEGMENT_BEGIN id=A12 source=colaboracion/parte_iii/source/public_final/ascensos_20260826/16f_realizacion_analitica_contraangulo.tex body_lines=574-748
\label{p3-20260826:sec:transportador-relativo-split}

La comparaison avec le bloc central d'APP peut s'effectuer sans identifier
des représentations distinctes. Fixons la base ordonnée des canaux et
l'involution
\[
 R=\begin{pmatrix}0&1\\1&0\end{pmatrix},
 \qquad R^2=I,
 \qquad
 P_\pm=\frac{I\pm R}{2}.
\]
Dans l'algèbre commutative
\[
 \mathscr B_R:=\{xI+yR:x,y\in\mathbb R\}
\]
considérons, dans la carte sexagésimale déjà déclarée,
\[
 B_0:=7I+2R,
 \qquad
 B_1:=\Adeg I+\Cdeg R.
\]
Le changement orthogonal de base
\[
 S=\frac1{\sqrt2}
 \begin{pmatrix}1&1\\1&-1\end{pmatrix}
\]
diagonalise simultanément l'involution et les deux blocs. Pour \(x>|y|\),
définissons
\[
 \rho(x,y):=\sqrt{x^2-y^2},
 \qquad
 \eta(x,y):=\operatorname{artanh}\!\left(\frac yx\right).
\]
Alors
\[
 xI+yR=\rho(x,y)e^{\eta(x,y)R}.
\]
Si
\[
 \rho_0=\sqrt{45},\quad
 \eta_0=\operatorname{artanh}(2/7),\quad
 \rho_1=\rho(\Adeg,\Cdeg),\quad
 \eta_1=\eta(\Adeg,\Cdeg),
\]
l'unique multiplication à gauche qui envoie \(B_0\) sur \(B_1\) est
\[
 \boxed{
 T_{\rm rel}
 :=\frac{\rho_1}{\sqrt{45}}
 e^{(\eta_1-\eta_0)R},
 \qquad
 T_{\rm rel}B_0=B_1.}
\]

\begin{proposition}[Transporteur relatif des deux blocs]
\label{p3-20260826:prop:transportador-relativo-split}
Une fois fixés \(R\), l'ordre des canaux, la branche positive et la carte
sexagésimale, \(T_{\rm rel}\) est l'unique élément de
\(\mathscr B_R\) qui satisfait \(TB_0=B_1\). Dans la base spectrale, il agit par
\[
 9\longmapsto\Adeg+\Cdeg,
 \qquad
 5\longmapsto\Adeg-\Cdeg.
\]
\end{proposition}

\begin{proof}
La décomposition polaire scindée donne la formule imprimée et prouve l'existence.
Comme \(\det B_0=45\ne0\), toute solution matricielle de \(TB_0=B_1\) satisfait
\(T=B_1B_0^{-1}\) ; elle est donc unique. La diagonalisation au moyen de \(S\)
produit les deux quotients spectraux
\((\Adeg+\Cdeg)/9\) et \((\Adeg-\Cdeg)/5\), équivalents à l'expression
exponentielle de \(T_{\rm rel}\).
\end{proof}

Ce transporteur n'est pas un entrelaceur représentationnel. En effet, un
entrelaceur \(\Phi\) entre les endomorphismes \(B_0\) et \(B_1\) devrait
satisfaire
\[
 \Phi B_0=B_1\Phi.
\]

\begin{theorem}[Obstruction spectrale à l'entrelacement]
\label{p3-20260826:thm:no-go-entrelazador-bloques-split}
Pour les valeurs de \(\Adeg\) et \(\Cdeg\) construites dans cette section,
\[
 \boxed{\Phi B_0=B_1\Phi\quad\Longrightarrow\quad\Phi=0.}
\]
\end{theorem}

\begin{proof}
Les spectres sont
\[
 \operatorname{Spec}(B_0)=\{9,5\},
 \qquad
 \operatorname{Spec}(B_1)
 =\{\Adeg+\Cdeg,\Adeg-\Cdeg\}.
\]
Les bornes rationnelles préalablement démontrées fournissent les intervalles
disjoints
\[
 7.29<\Adeg<7.30,
 \qquad
 2.12<\Cdeg<2.13,
\]
et, par conséquent,
\[
 9.41<\Adeg+\Cdeg<9.43,
 \qquad
 5.16<\Adeg-\Cdeg<5.18.
\]
Dans la base qui diagonalise \(R\), chaque entrée \(\phi_{ij}\) de \(\Phi\)
est multipliée par la différence entre une valeur propre de \(B_0\) et une de
\(B_1\). Toutes ces différences sont non nulles ; donc toutes les entrées de
\(\Phi\) s'annulent.
\end{proof}

Numériquement,
\[
 \rho_1=6.9814819335172378048\ldots,
 \qquad
 \frac{\rho_1}{\sqrt{45}}=1.0407378791354140779\ldots,
\]
\[
 \eta_1-\eta_0
 =0.005796351470571295590\ldots.
\]
La formule conserve l'algèbre scindée, les projecteurs \(P_\pm\) et la forme
lorentzienne à un facteur conforme \(\rho_1^2/45\) près. Elle dépend cependant
de \(B_1\) déjà construit : elle ne constitue pas une dérivation antérieure de
\(\Adeg\) ou \(\Cdeg\) et ne doit pas être confondue avec les entrelaceurs
équivariants entre réalisations d'un même module.

Pour
\[
S_{90}(q)=\frac{q^{90}}{1-q^{270}},
\qquad
S_{120}(q)=\frac{q^{120}}{1-q^{360}},
\]
la fonctionnelle de transition hexade--octade s'écrit intégralement dans la
carte des phases :
\[
\Gamma_{6\to8}
=\frac{180}{\pi}\Arad\Crad
+12\bigl(S_{90}(q_-)-S_{90}(q_+)\bigr)
-\bigl(S_{120}(q_-)-S_{120}(q_+)\bigr).
\]
L'évaluation produite par le contre-angle du
\cref{p3-20260826:thm:contraangulo-regional} est
\[
\boxed{
\Gamma_{6\to8}
=0.274007672694459274747255260774\ldots.}
\]
Les formes équivalentes de la fonctionnelle mathématique interne sont
\[
\boxed{
\Gamma_{6\to8}
=\Arad\Cdeg+K_{6\to8}
=\frac{180}{\pi}\Arad\Crad+K_{6\to8}
=\frac{\pi\Adeg\Cdeg}{180}+K_{6\to8}.}
\]
Ces identités évaluent \(\Gamma_{6\to8}\) après la génération de
\((\Crad,\Cdeg)\) ; elles ne sont pas utilisées en sens inverse pour sélectionner le
contre-angle. La lecture physique dans la connexion d'Ashtekar--Barbero est
représentée ensuite au moyen de
\[
 \Gamma_{6\to8}\dashrightarrow\gamma_{\rm BI}.
\]
La flèche n'introduit pas de troisième fonctionnelle et ne rétroagit pas sur \(C^*\). Les
séries \(S_{90}\), \(S_{120}\), les coefficients \(12\) et \(-1\) et leur
affectation au pas \(6\to8\) appartiennent à la construction mathématique
déclarée. L'identification physique exige en outre la normalisation propre
de la connexion.

% HMT_SOURCE_SEGMENT_END id=A12
\section{Naturalité sous raffinement et séparation d'avec la reconnaissance métrologique}
\label{sec:s102:c38:7}

\subsection{Estimation exacte du reste de la prolongation analytique}
\label{sec:s102:c38:7:1}

% HMT_SOURCE_SEGMENT_BEGIN id=B09 source=colaboracion/parte_iii/source/public_final/base_83/c30_body.tex body_lines=474-499

L'expression initialement trouvée,
\[
y_7=\frac\pi{90}(H_5+\alpha)-169\alpha^6-D_A\alpha^7,
\]
est la troncature de degré sept du caractère complet. Elle produit
\[
C_7=2.12373833896864600265403827103919\ldots{}^\circ.
\]
Le reste omis n'est pas un terme indéterminé ; c'est la queue géométrique
exacte
\[
\mathscr C_\pi(\alpha)
-169\alpha^6-D_A\alpha^7
=\frac{D_A^2\alpha^8}{1-D_A\alpha}.
\]
Par conséquent,
\[
\left|C_7-C_*^{\rm HMT}\right|
=\frac{180}{\pi}\frac{D_A^2\alpha^8}{1-D_A\alpha}
=2.7839208689064583\ldots\times10^{-16}\;{}^\circ.
\]
La récurrence transforme ainsi une approximation de septième degré en une
réalisation analytique complète et fournit une borne close pour chaque
troncature ultérieure.

% HMT_SOURCE_SEGMENT_END id=B09
\subsection{Dépendance fonctionnelle du sélecteur régional et du prolongement analytique}
\label{sec:s102:c38:7:2}

% HMT_SOURCE_SEGMENT_BEGIN id=B10 source=colaboracion/parte_iii/source/public_final/base_83/c30_body.tex body_lines=501-521

La construction change lorsque l'une quelconque de ses composantes actives est retirée.
Les valeurs suivantes s'obtiennent sans réétalonner les autres termes :
\[
\begin{array}{@{}lc@{}}
\toprule
\text{variante}&\text{coordonnée angulaire conjuguée en degrés}\\
\midrule
\rho_\pi=168&2.1237383389772976863904487\ldots\\
\rho_\pi=169&2.1237383389686457242619514\ldots\\
\rho_\pi=170&2.1237383389599937621334541\ldots\\
\text{sans queue déterminantale}&2.1237383389686949415301675\ldots\\
D_A\text{ remplacé par }1&2.1237383389686313409965826\ldots\\
\bottomrule
\end{array}
\]
Ces ablations ne constituent pas à elles seules une estimation probabiliste.
Elles démontrent que le sélecteur équivariant et le caractère déterminantal sont
des composantes effectives de la sortie et qu'ils ne peuvent être supprimés tout en maintenant la
même réalisation.

% HMT_SOURCE_SEGMENT_END id=B10
\subsection{Interprétation opératoire de la coordonnée angulaire conjuguée}
\label{sec:s102:c38:7:3}

% HMT_SOURCE_SEGMENT_BEGIN id=B14 source=colaboracion/parte_iii/source/public_final/base_83/c30_body.tex body_lines=677-715

La procédure complète peut être résumée par une chaîne d'applications :
\[
\begin{aligned}
\mathcal F_\pi
&\longmapsto b_\pi
\longmapsto\rho_\pi=169,\\
\alpha
&\longmapsto A
\longmapsto x
\longmapsto D_A,\\
(\rho_\pi,D_A)
&\longmapsto\mathscr C_\pi(\alpha),\\
(H_5,\mathscr C_\pi(\alpha))
&\longmapsto y_*
\longmapsto C_*^{\rm HMT},\\
(x,y_*)
&\longmapsto(q_+,q_-)
\longmapsto\Gamma_{6\to8}.
\end{aligned}
\]
La première ligne appartient au catalogue fini ; la deuxième, à la clôture commune de
\(\alpha\) ; la troisième, à la prolongation déterminantale ; la quatrième produit la
coordonnée orientée ; la cinquième la transporte vers l'interface dimensionnelle.

Dans cette organisation, \(\alpha\) mesure la contraction commune du transport et
\(C_*^{\rm HMT}\) mesure la séparation orientée de ses deux feuillets. Le retour
régional n'ajoute pas de constante externe : il sélectionne une amplitude finie et la
prolonge au moyen du déterminant déjà fixé par \(\alpha\). La carte en
degrés est ultérieure. L'objet primaire est le couple de phases \((x,y_*)\).

Une vérification exhaustive reconstruit les cinq régions directement
à partir du catalogue, parcourt les actions de
\(\mathfrak S_5\) et \(\mathfrak S_3\), vérifie la récurrence et somme la
série. L'entrée de \(\alpha\) est la sortie exacte de sa clôture
dodécaphasique, non un chiffre métrologique transcrit. La vérification maintient
séparés les théorèmes du catalogue, les règles du lecteur et la comparaison
externe ; c'est pourquoi aucune comparaison expérimentale n'intervient dans la fonction
génératrice.
% HMT_SOURCE_SEGMENT_END id=B14
\subsection{Certification de convergence de la troncature parangulaire}
\label{sec:s102:c38:7:4}

% HMT_SOURCE_SEGMENT_BEGIN id=A09 source=colaboracion/parte_iii/source/public_final/ascensos_20260826/16f_realizacion_analitica_contraangulo.tex body_lines=505-530

L'expression initialement trouvée,
\[
y_7=\frac\pi{90}(H_5+\alpha)-169\alpha^6-D_A\alpha^7,
\]
est la troncature de degré sept du caractère complet. Elle produit
\[
C_7=2.12373833896864600265403827103919\ldots{}^\circ.
\]
Le reste omis n'est pas un terme indéterminé ; c'est la queue géométrique
exacte
\[
\mathscr C_\pi(\alpha)
-169\alpha^6-D_A\alpha^7
=\frac{D_A^2\alpha^8}{1-D_A\alpha}.
\]
Par conséquent,
\[
\left|C_7-\Cdeg\right|
=\frac{180}{\pi}\frac{D_A^2\alpha^8}{1-D_A\alpha}
=2.7839208689064583\ldots\times10^{-16}\;{}^\circ.
\]
La récurrence transforme ainsi une approximation de septième degré en une
réalisation analytique complète et fournit une borne close pour chaque
troncature ultérieure.

% HMT_SOURCE_SEGMENT_END id=A09
\subsection{Naturalité du sélecteur régional sous raffinement et changement de carte}
\label{sec:s102:c38:7:5}

% HMT_SOURCE_SEGMENT_BEGIN id=A10 source=colaboracion/parte_iii/source/public_final/ascensos_20260826/16f_realizacion_analitica_contraangulo.tex body_lines=532-552

La construction change lorsque l'une quelconque de ses composantes actives est retirée.
Les valeurs suivantes s'obtiennent sans réétalonner les autres termes :
\[
\begin{array}{@{}lc@{}}
\toprule
\text{variante}&\text{contre-angle en degrés}\\
\midrule
\rho_\pi=168&2.1237383389772976863904487\ldots\\
\rho_\pi=169&2.1237383389686457242619514\ldots\\
\rho_\pi=170&2.1237383389599937621334541\ldots\\
\text{sans queue déterminantale}&2.1237383389686949415301675\ldots\\
D_A\text{ remplacé par }1&2.1237383389686313409965826\ldots\\
\bottomrule
\end{array}
\]
Ces ablations ne constituent pas à elles seules une estimation probabiliste.
Elles démontrent que le sélecteur équivariant et le caractère déterminantal sont
des composantes effectives de la sortie et qu'ils ne peuvent être supprimés tout en maintenant la
même réalisation.

% HMT_SOURCE_SEGMENT_END id=A10
\subsection{Indépendance métrologique de la construction parangulaire}
\label{sec:s102:c38:7:6}

% HMT_SOURCE_SEGMENT_BEGIN id=A14 source=colaboracion/parte_iii/source/public_final/ascensos_20260826/16f_realizacion_analitica_contraangulo.tex body_lines=799-831

La procédure complète peut être résumée par une chaîne d'applications :
\[
\begin{aligned}
\mathcal F_\pi
&\longmapsto b_\pi
\longmapsto\rho_\pi=169,\\
\alpha
&\longmapsto \Adeg
\longmapsto \Arad
\longmapsto D_A,\\
(\rho_\pi,D_A)
&\longmapsto\mathscr C_\pi(\alpha),\\
(H_5,\mathscr C_\pi(\alpha))
&\longmapsto \Crad
\longmapsto \Cdeg,\\
(\Arad,\Crad)
&\longmapsto(q_+,q_-)
\longmapsto\Gamma_{6\to8}.
\end{aligned}
\]
La première ligne appartient au catalogue fini ; la deuxième, à la clôture commune de
\(\alpha\) ; la troisième, à la prolongation déterminantale ; la quatrième produit la
coordonnée orientée ; la cinquième la transporte vers l'interface dimensionnelle.

Dans cette organisation, \(\alpha\) mesure la contraction commune du transport et
le contre-angle mesure la séparation orientée de ses deux feuillets. Dans
l'évaluation en vigueur, le retour régional sélectionne une amplitude finie et la
prolonge au moyen du déterminant déjà fixé par \(\alpha\), sans insérer de
constante externe supplémentaire. La carte en degrés est ultérieure. L'objet
premier est le couple de phases \((\Arad,\Crad)\). La récurrence, la somme
régionale et les ablations dépendent du catalogue et de
\(\alpha_{\rm HMT}\), mais ne reçoivent pas de constante métrologique supplémentaire.
% HMT_SOURCE_SEGMENT_END id=A14

